\section{Anexos: Documentación del Código}

En esta sección se presenta el código fuente de los scripts desarrollados en Python para el análisis numérico del farmeo. Cada script incluye una explicación de su propósito y la función que cumple dentro del proyecto.

\subsection{Script \texttt{fopt.py}: Cálculo del Farmeo Óptimo}
\textbf{Propósito:} Define las funciones matemáticas para calcular el oro teórico de los creeps y compara los resultados del modelo con los datos reales de jugadores Inmortales y Promedio.

\begin{lstlisting}[language=Python, frame=single, rulecolor=\color{black}, backgroundcolor=\color{backcolour}, caption=Código principal de \texttt{fopt.py}]
import numpy as np
import matplotlib.pyplot as plt
import os

def c_melee(t):
    return 3 + (t >= 15) + (t >= 30) + (t >= 45)

def c_ranged(t):
    return 1 + (t >= 40)

def g_oleada(t):
    return 36 * c_melee(t) + 48 * c_ranged(t)

def f_opt(T, r=0.5, g_camp=78):
    # Suma de oleadas y campamentos
    oleadas = sum(g_oleada(k / 2) for k in range(2 * T))
    neutrales = r * g_camp * T
    return oleadas + neutrales
\end{lstlisting}

\subsection{Script \texttt{interpolacion.py}: Interpolación de Lagrange}
\textbf{Propósito:} Implementa el método de interpolación de Lagrange para estimar el oro acumulado en minutos no medidos, a partir de datos discretos conocidos.

\begin{lstlisting}[language=Python, frame=single, rulecolor=\color{black}, backgroundcolor=\color{backcolour}, caption=Código principal de \texttt{interpolacion.py}]
def lagrange(x_puntos, y_puntos, x):
    total = 0.0
    n = len(x_puntos)
    for i in range(n):
        term = y_puntos[i]
        for j in range(n):
            if i != j:
                # Calculo del termino de Lagrange
                term *= (x - x_puntos[j]) / (x_puntos[i] - 
                         x_puntos[j])
        total += term
    return total
\end{lstlisting}

\subsection{Script \texttt{error\_relativo.py}: Análisis del Error}
\textbf{Propósito:} Calcula el porcentaje de error relativo entre el farmeo óptimo (Inmortal) y el farmeo real (Promedio) para los tres héroes, y genera una gráfica de su evolución temporal.

\begin{lstlisting}[language=Python, frame=single, rulecolor=\color{black}, backgroundcolor=\color{backcolour}, caption=Código principal de \texttt{error\_relativo.py}]
import pandas as pd

# Lectura de datos reales
df_real = pd.read_csv('../datos/farmeo_real.csv')

for heroe in df_real['heroe'].unique():
    sub = df_real[df_real['heroe'] == heroe].sort_values('minuto')
    for _, row in sub.iterrows():
        T = int(row['minuto'])
        Fo = f_opt(T)
        Fr = row['oro_real']
        # Calculo del error relativo
        Er = abs(Fo - Fr) / Fo * 100 
\end{lstlisting}

\subsection{Script \texttt{grafo\_rutas.py}: Modelado de Rutas}
\textbf{Propósito:} Construye un grafo dirigido para representar los campamentos neutrales y la línea del mapa, asignando pesos a las aristas según el tiempo de desplazamiento, con el fin de optimizar la ruta de farmeo.

\begin{lstlisting}[language=Python, frame=single, rulecolor=\color{black}, backgroundcolor=\color{backcolour}, caption=Código principal de \texttt{grafo\_rutas.py}]
import networkx as nx
import matplotlib.pyplot as plt

# Creacion del grafo dirigido
G = nx.DiGraph()
G.add_edge('linea', 'campA', tiempo=14, oro=78)
G.add_edge('campA', 'linea', tiempo=14, oro=0)
G.add_edge('linea', 'campB', tiempo=20, oro=78)
\end{lstlisting}